\section{Discussion of the Role of Joint Prediction}
\label{app:joint-prediction}

The conditional-risk analysis describes how action supervision shapes representations; the planning discussion connects learned proposals to contact-dependent outcomes. These mechanisms complement Sections~\ref{sec:what-joint-buys} and~\ref{sec:control-planning}.

\Needspace{30\baselineskip}
\phantomsection
\label{app:system-simplicity}

Table~\ref{tab:world-action-factorizations} complements Section~\ref{sec:related} with a
comparison of training requirements and prediction capabilities.

% Evidence: original papers linked by each method's citation.
% Local audit: local/literature/paradigm-table-2026-09-13/README.md.
% Joint means clean-action FDM and policy losses in the SAME training pass.
% It does not mean generic joint video/action generation, which several baselines support.
\begin{table}[H]
\centering
\caption{\lead{Modeling and Control Capabilities.}
FDM denotes forward dynamics conditioned on clean actions; Policy denotes learned action
prediction. Entries describe reported training recipes or the goal-reaching configurations used here.}
\label{tab:world-action-factorizations}
\begingroup
\newcommand{\yes}{\ding{51}}
\newcommand{\no}{\ding{55}}
\setlength{\tabcolsep}{3pt}
\renewcommand{\arraystretch}{1.05}
\begin{tabular*}{\linewidth}{@{\extracolsep{\fill}}lcccc@{}}
\toprule
Method & FDM & Policy & \shortstack{Joint FDM\\and Policy} & \shortstack{End-to-end\\from scratch} \\
\midrule
\multicolumn{5}{@{}l}{\emph{World models}} \\
\methodname{DINO-WM}~\citep{zhou2024dinowm} & \yes & \no & \no & \no \\
\methodname{V-JEPA 2-AC}~\citep{assran2025vjepa2} & \yes & \no & \no & \no \\
\methodname{LeWM (CEM)}~\citep{maes2026lewm} & \yes & \no & \no & \yes \\
\midrule
\multicolumn{5}{@{}l}{\emph{World action models}} \\
\methodname{UVA}~\citep{li2025uva} & \yes & \yes & \no & \no \\
\methodname{UWM}~\citep{zhu2025uwm} & \yes & \yes & \no & \no \\
\methodname{DreamZero}~\citep{ye2026dreamzero} & \no & \yes & \no & \no \\
\methodname{Fast-WAM}~\citep{yuan2026fastwam} & \no & \yes & \no & \no \\
\methodname{Cosmos 3}~\citep{nvidia2026cosmos3} & \yes & \yes & \no & \no \\
\textbf{\method{} (ours)} & \yes & \yes & \yes & \yes \\
\midrule
\multicolumn{5}{@{}l}{\emph{Goal-conditioned policy}} \\
\methodname{LeWM (IDM)}~\citep{nguyen2026gcidm} & \yes & \yes & \no & \no \\
\methodname{GCBC}~\citep{ghosh2021gcsl} & \no & \yes & \no & \no \\
\bottomrule
\end{tabular*}
\par\vspace{3pt}
\begin{minipage}{\linewidth}
Joint requires clean-action FDM and policy losses on the same sample in one training pass.
From scratch includes visual encoding and prediction without pretrained or frozen components.
UWM uses a frozen pretrained variational autoencoder even when its robot backbone is trained from scratch.
\end{minipage}
\endgroup
\end{table}


\Needspace{12\baselineskip}
\subsection{Representation and Dynamics}
\label{app:joint-representation}

\begin{wrapfigure}{r}{0.40\linewidth}
\vspace{-16pt}
\centering
\includegraphics[width=\linewidth,trim=20bp 25bp 20bp 30bp,clip]{figures/assets/anti_collapse/anti_collapse_open_arrows.pdf}
\setlength{\abovecaptionskip}{3pt}
\vspace{-4pt}
\caption{Conceptual roles of SIGReg and action flow matching in shaping the latent space.}
\label{fig:anti-collapse-mechanisms}
\vspace{-6pt}
\end{wrapfigure}

\lead{Representation Shaping.} SIGReg~\citep{balestriero2025lejepa,maes2026lewm}
regularizes the latent marginal; action supervision encourages control-relevant
distinctions (Figure~\ref{fig:anti-collapse-mechanisms}). If visual history
changes the conditional action velocity, discarding it increases optimal action
regression risk. Latent prediction adds pressure to retain action-conditioned
consequences.
\finishwrapfigure

\Needspace{10\baselineskip}
\lead{Conditional Action Risk.} In Equation~\ref{eq:action-flow-path}, squared
flow matching makes the optimal velocity the conditional mean of
$\va_j-\vepsilon_j$. Let $\mathcal I_j$ include the noisy action prefix, flow
time, and other conditioning, but exclude visual history $\vz_{\leq t}$. With
$\mathcal R_{\mathrm{FM}}$ denoting optimal velocity MSE, the conditional
projection identity gives
\begin{equation}
    \begin{aligned}
    &\mathcal{R}_{\mathrm{FM}}(\mathrm{constant})-
      \mathcal{R}_{\mathrm{FM}}(\vz_{\leq t})\\
    &=\frac{1}{C}\sum_{j=t}^{t+C-1}\mathbb{E}\!\left[
      \operatorname{MSE}\!\left(
      \mathbb{E}[\va_j-\vepsilon_j\mid\mathcal I_j,\vz_{\leq t}],
      \mathbb{E}[\va_j-\vepsilon_j\mid\mathcal I_j]
      \right)
    \right]\geq 0.
    \end{aligned}
    \label{eq:flow-collapse-gap}
\end{equation}
When visual history changes the conditional velocity, a constant code incurs higher optimal risk. Action supervision therefore favors these distinctions, while SIGReg regularizes the global latent distribution.

\Needspace{24\baselineskip}
\subsection{Planning with Learned Action Proposals}
\label{app:joint-planning}

\begin{figure}[H]
\centering
\includegraphics[width=\linewidth]{figures/contact_geometry.pdf}
\caption{\lead{Action Proposals and Predicted Consequences.}
\textbf{(a)} Smooth rollouts; \textbf{(b)} candidate rollouts branch at contact;
\textbf{(c)} LeWAM pairs learned action proposals with predicted outcomes.
Dashed links denote search; the solid link denotes joint prediction.
}
\label{fig:contact-geometry}
\end{figure}

\lead{Ill-Posed Action Solver in Contact-rich Events.} Prior JEPAs for control \eg, DINO-WM ~\citep{zhou2024dinowm} and \methodname{LeWM}~\citep{maes2026lewm}
learn action-conditioned latent dynamics without an action-prediction
objective and uses CEM~\citep{rubinstein1999cem} for downstream control. In smooth dynamics like 2D navigation, the action solver is \emph{well-posed}: a diverse representation can invert the action trajectories by argmin the cost of action-conditioned dynamics. However, a contact event can switch the active mode, so
nearby actions yield distinct outcomes~\citep{pang2023contactplanning}. If
sampled plans miss useful contact, an object-based goal cost provides little
ranking signal~\citep{yamada2024dcubed}. With a finite population, CEM ranks the contact branches it samples, and additional iterations incur more rollouts. 


\lead{Behavior and Consequence.}  Joint prediction
additionally supplies a learned action distribution for planning proposals.
Learned proposals can improve contact-rich planning under
constrained budgets~\citep{brudermuller2025gpc}. Figure~\ref{fig:contact-geometry}
illustrates why the population-based search from a broad prior cannot work on contact-rich manipulation, and how LeWAM uses the learned proposals and predicted consequences to alleviate this problem. The
planning objective is defined in Section~\ref{sec:control-planning}, and the
refinement procedure is specified in Appendix~\ref{app:planners}. Though this problem in 3D contact-rich manipulation is well known and our solution is heuristic, we consider the rigorous verification of the claim as future work.


\clearpage
\section{Author Contribution Statement}
\label{app:author-contributions}
All coauthors contributed to project discussions, writing, and editing.

Minghao initiated the project and led its technical development. He proposed the
general architecture with action flow matching, MoT, the initial version of
gradient-based planning, and constructed the benchmarks. He also wrote the main
part of the paper and built the webpage.

Tavis designed the details of model architecture and MoT attention mechanism, optimized the planning algorithm to use batched gradients, and refined the methodology to the final version through extensive experimentation. He also finalized the released code. Tavis and Minghao
completed the simulation experiments reported in the paper.

Eryk designed and carried out the entire physical robot study. He built the hardware and teleoperation pipeline, adapted the model for real-robot deployment, collected all real-robot data, trained models for real robots, and ran the physical evaluations. He also wrote the main part of physical-robot experiments section and prepared its results.

Randall initiated and supervised the project and guided its overall direction. He advised on representation learning, experimental design, and result verification, contributed extensively to writing and refining the manuscript, and finalized the paper's title, narrative, and branding.
